\chapter{Interaction}

\section{From Intent to Action}

Chapter 10 gave affordance a precise meaning: at the meeting of field and capability, a menu of candidate actions, P(\(\Delta\psi\)), becomes available to an agent. Chapter 13 gave the machinery that turns a selected candidate into an actual proposal, combining the field's intrinsic tendencies with a learned forcing term and choosing among live candidates by relevance. Something is still missing between these two accounts, and it is the subject of this chapter: neither Chapter 10 nor Chapter 13 ever specified how an agent's actual intent, a player's click, a keypress, another AI system's planned action, enters this pipeline in the first place. This chapter builds the front end that both of those chapters assumed already existed.

\section{The Action Interface}

Rather than exposing an open-ended, unstructured action space, this architecture's interaction layer is built from the small operational vocabulary Chapter 10 already introduced and deliberately kept narrow: store, retrieve, modify, compose, and rewrite. Every action any agent takes, however complex it appears from the outside, resolves into some combination of these five primitives applied to some target patch or region. Picking up an object is retrieval followed by a compose operation binding it to the agent's own state. Placing a bench is a store. Repairing a crack is a modify. Combining two ingredients into a meal is a compose. Rewriting a sign's text according to a rule, rather than an arbitrary edit, is, unsurprisingly, a rewrite. This is a genuine constraint on what the interaction layer will accept, not merely a convenient way of describing actions after the fact, and the constraint is worth keeping: an interaction layer that accepted arbitrary, unstructured intent would have no principled way to check whether a given action even has the shape of something the affordance and proposal machinery downstream can process. Five operations are enough to cover the space this book cares about, and narrowing the interface to exactly these five is what keeps the rest of this chapter's architecture tractable.

\section{A Worked Example: Two Orders for One Alphabet}

It is worth working through a single example in enough depth to show what the rewrite operation from section 28.2 is actually for, because "replace content according to some rule rather than an arbitrary edit" can sound abstract until it is seen doing real work on something concrete.

The Arabic alphabet's twenty-eight letters have, over their history, been given two distinct total orderings. The Abjad order is the older of the two, inherited directly from the Phoenician and Aramaic sequences that preceded it, and it is not arbitrary: it doubles as a numbering system, Abjad numerals, in which each letter also stands for a specific number, used historically in chronology, astronomy, and other contexts where a script needed to double as a counting system. The order is preserved today chiefly through mnemonic strings, Abjad, Hawwaz, Huṭṭī, Kalaman, Saʿfaṣ, Qarashat, Thakhadh, Ḍaẓagh, that let the sequence be memorized as if it were itself a set of words. The Hijā'ī order is considerably younger, adopted roughly a century after the start of Islam, and it reorganizes the same twenty-eight letters around an entirely different principle: letters that share a common base shape, distinguished from one another only by the placement of diacritical dots, are grouped adjacently, which makes the sequence considerably easier to teach to someone learning to write, and it is this order, not Abjad, that modern dictionaries, indices, and directories actually use.

Two features of this history matter for the interaction layer this chapter has been building. First, the transition from Abjad to Hijā'ī is a rewrite in exactly this book's technical sense, not a modify and not a wholesale replacement. Every one of the twenty-eight letters survives the transition; nothing is added, and nothing is deleted. What changes is a relation, the adjacency structure imposed on the letters, according to an explicit, statable rule: letters with the same base shape belong together. This is precisely the distinction section 28.2 drew between rewrite and modify, and precisely why this book insists on keeping rewrite as its own primitive rather than treating it as a large modify: a rewrite is content-preserving and rule-governed in a way an arbitrary edit is not, and the Abjad-to-Hijā'ī transition is a case where every letter's persistence through the change can be checked, the way this book's own conservation laws, from Chapter 22, would check that nothing was silently lost in a large-scope proposal.

Second, and just as important, adopting Hijā'ī order for dictionaries and instruction did not retire Abjad order. Abjad numerals are still the ordering in active use wherever the numerical reading of the letters actually matters, and nothing about Hijā'ī's dominance in one context makes Abjad wrong in another. This is Chapter 10's affordance argument in a different register: what structure a set of letters ought to have is not an intrinsic fact about the letters, it depends on what the structure is being asked to do, exactly as a ladder's affordance depends on who is climbing it rather than on the ladder alone. A single alphabet, asked to serve numerology, asks for one order; the same alphabet, asked to serve handwriting pedagogy, asks for another; and this architecture has no trouble holding both answers at once, because neither is more true of the letters than the other.

If this book's write cycle had to process the historical adoption of Hijā'ī order as a single proposal, it is worth noting what would make that proposal admissible rather than simply large. A \(\Delta\psi\) that reordered the alphabet while silently dropping a letter, or introducing one that was never there, would fail Chapter 22's conservation requirements outright, whatever pedagogical convenience it offered. A \(\Delta\psi\) that reordered the alphabet according to no statable rule at all, an arbitrary shuffle rather than a shape-based regrouping, would still conserve the letters but would no longer be a rewrite in the sense this chapter's interface actually accepts; it would be indistinguishable, from the interaction layer's point of view, from vandalism. What made Hijā'ī's adoption a legitimate rewrite, worth committing rather than refusing, was precisely that it was both content-preserving and rule-governed, the two properties this book has insisted on for every large-scope proposal since section 28.2 first drew the line between rewrite and arbitrary edit.

\section{Cues as the Bridge from Perception to Action}

The full path from perception to action now has a shape worth tracing end to end. An agent's observation, arriving through Chapter 27's rendering path, surfaces cues in the sense Chapter 10 developed: signals that activate relevance, \(\rho_c\), for whatever capabilities the agent brings. Those activated cues make some subset of the local affordance menu live, exactly as Chapter 10 described. Somewhere outside this book's scope, in the agent's own cognition or planning process, whether that agent is a human mind or another AI system, a selection is made among the live affordances. This chapter's interaction layer receives that selection, expressed in terms of the five-operation vocabulary from section 28.2, and hands it forward as the a that Chapter 13's selection policy, π(a|x) ∝ exp(β·\(\rho_c(x,a)\)), was already built to receive. Nothing in this chapter tries to explain how an agent decides what it wants. That decision is the agent's own, made by whatever cognitive or computational process constitutes it. What this chapter supplies is the interface an already-made decision passes through on its way into the world's own write cycle.

\section{Human and AI Agents, Same Interface}

It is worth stating explicitly something the last section only implied: this architecture does not distinguish, at the interaction layer, between a human agent and an AI agent. A human player's click or keypress and another AI system's action output are, from the interface's own point of view, the same kind of thing, a selection among live affordances expressed in the five-operation vocabulary. Both arrive at the proposal service, per Chapter 25's architecture, exactly the same way, and both are held to exactly the same admissibility discipline once Chapter 14's projection stage receives them. This uniformity is not an incidental simplification. It is what makes the persistent NPCs and multi-agent applications Chapter 31 will eventually build possible at all: an NPC's actions and a human player's actions are proposals of the same kind, checked against the same constraint manifold, capable of conflicting with or complementing one another exactly as any two simultaneous proposals would, per Chapter 15's arbitration.

\section{Direct Editing Is Still a Proposal}

It is tempting to imagine a privileged path for tools that need to edit the world directly, a level designer's editor, an administrative console, a debugging interface, bypassing the affordance-cue-selection pipeline sections 28.1, 28.2, and 28.4 developed, on the grounds that such tools are not really agents acting within the world but external instruments operating on it. This chapter argues against any such bypass, and the reasoning is the same reasoning Chapter 27.4 already established for rendering-triggered invention: a write that skips \(\Pi_C\) and Commit is a write with no guarantee that it respects the entropy constraint, appearance consistency, or accumulated residue, and an editor able to silently violate these is an editor able to introduce exactly the kind of unearned certainty and unaccountable contradiction this entire book has been built to prevent. A level designer's edit is still, architecturally, a \(\Delta\psi\). It may be a \(\Delta\psi\) with unusually broad scope, section 28.3's worked example was, in effect, exactly this kind of large-scope rewrite, and it may be submitted through a different front end than section 28.2's five-operation vocabulary, an editor plausibly needs a richer authoring surface than an in-world agent does, but it still enters the same proposal-projection-commit pipeline as any other proposed change, and it is still subject to being corrected or refused if it conflicts with what the field already holds. This gives the architecture a clean, memorable invariant worth stating on its own: every write to \(M_t\) is a \(\Delta\psi\), checked and committed through the same pipeline, regardless of what proposed it. There is no other way to change the world, and no component anywhere in this book's architecture holds an unchecked path to the world database that bypasses this rule.

\section{Looking Ahead}

This chapter has completed the picture for a single agent: intent, expressed through a narrow action vocabulary, entering the world's write cycle on exactly the same footing whether that agent is human, artificial, or an editing tool operating from outside the simulated world entirely. Nothing in this chapter has yet addressed what happens when many such agents act on the same world at overlapping moments, beyond the brief mention, in section 28.5, that their proposals are held to a common discipline. Chapter 29 takes up multiplicity directly, extending this chapter's single-agent interface to genuine multiplayer coordination.
